% !TEX TS-program = xelatex
% !BIB program = bibtex
% !TeX spellcheck = ru_RU
% !TEX root = talk.tex

\documentclass
  [ russian
  , aspectratio=43
  ] {beamer}


\input{preamble.tex}
\makeatletter

% ------------------------------------------

\usepackage[HTML, dvipsnames]{xcolor}
\setbeamercolor{item}{fg=red}

\usepackage{listings}
\usepackage{listingsutf8}

\lstdefinelanguage{ML}{
  keywords={type, match, with, let, rec, in, fun, function, of, if, then, else},
  keywordstyle=\color{Maroon}\bfseries,
  identifierstyle=\color{black},
  morecomment=[s]{(*}{*)},
  commentstyle=\color{gray},
  morestring=[b]",
  stringstyle=\color{green},
}

\lstset{
    basicstyle=\ttfamily\scriptsize,
    showstringspaces=false
}

\setmonofont{DejaVu Sans Mono}

\setbeamertemplate{itemize subitem}{\usebeamertemplate{itemize item}}
\setbeamertemplate{itemize subsubitem}{\usebeamertemplate{itemize item}}

\definecolor{codegreen}{rgb}{0,0.6,0}
\definecolor{codeblack}{rgb}{0,0,0}
\definecolor{colorprint1}{HTML}{E37DCB}
\definecolor{colorprint2}{HTML}{E3C87D}
\definecolor{colorprint3}{HTML}{7DE395}
\definecolor{colorprint4}{HTML}{7D98E3}

\usepackage{capt-of}
\usepackage{caption}
\captionsetup[figure]{labelformat=empty, font={scriptsize}}

% ------------------------------------------

\input{pretitle.tex}
\input{title.tex}

\newcommand{\academicGroup}{\my@title@group@ru}
\newcommand{\advisorChair}{\my@title@chair@ru}
\title[Самовоспроизводящийся компилятор]{\my@title@title@ru}
\author[\my@title@author@ru]{\my@title@author@ru, группа \academicGroup}
\institute[СПбГУ]{}
\date[2 июня 2026 г.]{}
\newcommand{\supervisor}{\my@title@supervisor@ru}
\newcommand{\supervisorPosition}{\my@title@supervisorPosition@ru}
\newcommand{\consultant}{\my@title@consultant@ru}
\newcommand{\consultantPosition}{\my@title@consultantPosition@ru}
\newcommand{\reviewer}{\my@title@reviewer@ru}
\newcommand{\reviewerPosition}{\my@title@reviewerPosition@ru}
\newcommand{\defenseYear}{\my@title@year@ru}

\makeatother
\begin{document}

\begin{frame}
    \includegraphics[width=1.4cm]{figures/SPbGU_Logo.png}
    \vspace{-35pt}
    \hspace{-10pt}
    \begin{center}
        \begin{tabular}{c}
            \scriptsize{Санкт-Петербургский государственный университет} \\
            \scriptsize{\advisorChair}
        \end{tabular}
        \titlepage
    \end{center}

    \btVFill

    {\scriptsize
        \textbf{Научный руководитель:}  \supervisorPosition~\supervisor \\
    }
    \makeatother
    \begin{center}
        \vspace{5pt}
        \scriptsize{Санкт-Петербург\\ \defenseYear}
    \end{center}
\end{frame}


\begin{frame}
    \frametitle{Постановка задачи}
    \textbf{Целью} работы является обеспечение возможности сборки языковых инструментов с помощью rukaml
    \vspace{1em}

    \textbf{Задачи}:
    \begin{itemize}
        \item Расширение стандартной библиотеки:
              \begin{itemize}
                  \item Работа с файловой системой
                  \item Обработка аргументов командной строки
                  \item Форматированный вывод
              \end{itemize}
        \item Реализация манглирования имён и разрешения внешних ссылок
        \item Реализация простого компилятора и его сборка с помощью rukaml
    \end{itemize}
\end{frame}


\begin{frame}
    \frametitle{Замечание}
    \begin{itemize}
        \item Все последующие примеры кода — программы, которые компилируются с помощью rukaml
        \item Эти примеры показывают, как написание кода выглядит для пользователя
        \item Они не имеют отношения к внутреннему устройству rukaml
    \end{itemize}
\end{frame}


\begin{frame}[fragile]
    \frametitle{Расширение стандартной библиотеки}

    Форматированный вывод

    \begin{lstlisting}[
        language=ML,
        columns=fullflexible,
        stringstyle=\color{codegreen},
        emph={fprintf},
        emphstyle=\bfseries\color{codeblack}
    ]
let rec pp_pexpr oc e =
    match e with
    | Pexpr_binop (op, e1, e2) ->
        fprintf oc "(%a %a %a)" pp_pexpr e1 pp_pbinop op pp_pexpr e2
    | Pexpr_assign (v, e) ->
        fprintf oc "%s = %a" v pp_pexpr e
    | Pexpr_tern (c, t, e) ->
        fprintf oc "%a ? %a : %a" pp_pexpr c pp_pexpr t pp_pexpr e
    | ...
    \end{lstlisting}

    Обработка аргументов командной строки и работа с файловой системой

    \begin{lstlisting}[
        language=ML,
        columns=fullflexible,
        emph={sys_argv, open_in, open_out, input_all, close_in, close_out},
        emphstyle=\bfseries\color{codeblack}
    ]
let main =
  let input_path, output_path = parse_args sys_argv in
  let input_channel = open_in input_path in
  let output_channel = open_out output_path in
  let () = run_program output_channel (input_all input_channel) in
  let () = close_in input_channel in
  let () = close_out output_channel in
  exit 0
    \end{lstlisting}

\end{frame}


\begin{frame}
    \frametitle{\fontsize{13pt}{15.6pt}\selectfont Манглирование имён и разрешение внешних ссылок [1/2]}
    \begin{itemize}
        \item Обеспечение валидности генерируемых ассемблерных меток и предотвращение их коллизий
        \item Поддержание таблиц для сопоставления ассемблерных меток исходным именам
        \item Разрешение синонимов для примитивов стандартной библиотеки
    \end{itemize}
\end{frame}


\begin{frame}[fragile]
    \frametitle{\fontsize{13pt}{15.6pt}\selectfont Манглирование имён и разрешение внешних ссылок [2/2]}
    Пример кода, для которого разрешение имён корректно работает

    \begin{lstlisting}[
    language=ML,
    escapechar=|,
    basicstyle=\ttfamily\footnotesize,
    columns=fullflexible
]
  (* |$\textcolor{colorprint1}{\textbf{print}}$| is defined somewhere |$\textcolor{gray}{\texttt{in}}$| stdlib *)

  let () = |$\textcolor{colorprint1}{\textbf{print}}$| 42
  let () =
      let |$\textcolor{colorprint2}{\textbf{print}}$| _ = () in
      |$\textcolor{colorprint2}{\textbf{print}}$| 42
  let () = |$\textcolor{colorprint1}{\textbf{print}}$| 42
  let |$\textcolor{colorprint3}{\textbf{print}}$| _ = ()
  let () = |$\textcolor{colorprint3}{\textbf{print}}$| 42
  let () =
      let |$\textcolor{colorprint4}{\textbf{print}}$| _ = () in
      |$\textcolor{colorprint4}{\textbf{print}}$| 42
  let () = |$\textcolor{colorprint3}{\textbf{print}}$| 42
\end{lstlisting}
\end{frame}


\begin{frame}[fragile]
    \frametitle{Компилятор, собираемый с помощью rukaml [1/3]}
    \begin{itemize}
        \item Создан с целью проверки возможностей rukaml
        \item Представляет собой простой компилятор подмножества языка C в язык ассемблера RISC-V 64
        \item Включает в себя:
              \begin{itemize}
                  \item Определение AST
                  \item Синтаксический анализатор
                  \item Pretty-printer
                  \item Модуль кодогенерации
                  \item Обработку аргументов командной строки, чтение входных данных, вывод результата работы
              \end{itemize}
    \end{itemize}
\end{frame}


\begin{frame}[fragile]
    \frametitle{Компилятор, собираемый с помощью rukaml [2/3]}

    Пример программы, которую он способен собрать

    \begin{lstlisting}[language=C, keywordstyle=\color{Maroon}\bfseries]
int fact_rec(int n) {
    if (n < 1) {
        return 1;
    }
    return n * fact_rec(n - 1);
}

int fact_iter(int n) {
    int acc = 1;
    for (int i = 2; i <= n; i = i + 1) {
        acc = acc * i;
    }
    return acc;
}

int main() {
    return 144 - fact_rec(4) - fact_iter(5);
}
\end{lstlisting}

\end{frame}


\begin{frame}
    \frametitle{Компилятор, собираемый с помощью rukaml [3/3]}

    \begin{columns}[c]
        \begin{column}{0.35\textwidth}
            \centering
            \includegraphics[width=\textwidth]{figures/scheme.png}
            \captionof{figure}{Упрощенная схема архитектуры}
        \end{column}
        \begin{column}{0.5\textwidth}
            \begin{itemize}
                \item \texttt{\scriptsize compiler} — компилятор, который собирается с помощью rukaml
                \item \texttt{\scriptsize program} — программа, которая собирается с помощью \texttt{\scriptsize compiler}
            \end{itemize}
        \end{column}
    \end{columns}
\end{frame}


\begin{frame}[fragile]
    \frametitle{Тестирование}
    \begin{itemize}
        \item В основном, писались системные тесты, проверяющие работу сразу всех этапов компиляции
        \item Компилятор, собираемый с помощью rukaml, тестировался путём сравнения поведения с ocamlopt
        \item Процент покрытия измерялся с помощью утилиты \texttt\small{bisect\_ppx}
              \begin{itemize}
                  \item Результат: 79~\% по строкам
              \end{itemize}
    \end{itemize}
\end{frame}


\begin{frame}[fragile]
    \frametitle{Результаты}
    \begin{itemize}
        \item Стандартная библиотека расширена ранее упомянутыми примитивами
        \item Реализовано манглирование имён и разрешение внешних ссылок
        \item Благодаря внесённым изменениям, с помощью rukaml удалось собрать простой компилятор
              \begin{itemize}
                  \item Принципиальных ограничений, не позволяющих собирать с помощью rukaml более сложные языковые инструменты, не выявлено
              \end{itemize}
        \item Выступление с докладом на секционном заседании конференции СПбПУ «Современные технологии в теории и практике программирования»
              \begin{itemize}
                  \item Тезисы доклада включены в сборник материалов конференции, проиндексированный в РИНЦ
              \end{itemize}
    \end{itemize}
\end{frame}


\end{document}
